\section{Methodology}

\subsection{System Overview}

ReAct SQL consists of two phases unified under the ReAct paradigm: (1) \textbf{Onboarding Phase} -- autonomous database analysis to build Rich Context, and (2) \textbf{Runtime Phase} -- tool-augmented SQL generation with self-correction. Both phases use the same ReAct loop (Thought-Action-Observation) but with different task objectives.

\subsection{Rich Context Construction}

\subsubsection{ReAct-Driven Database Analysis}

Unlike template-based approaches, we use task-driven database analysis. Given the task "analyze this database and collect all necessary information", the LLM agent autonomously decides which SQL queries to execute. A typical flow:

\begin{enumerate}
    \item \textbf{Thought}: "Need to identify database type and version"
    \item \textbf{Action}: \texttt{execute\_sql("SELECT version()")}
    \item \textbf{Observation}: "PostgreSQL 14.5"
    \item \textbf{Thought}: "Need to collect schema information"
    \item \textbf{Action}: \texttt{execute\_sql("SELECT ... FROM information\_schema.columns")}
    \item ... (continues iteratively)
\end{enumerate}

This approach adapts to different database types, scales, and features without manual configuration.

\subsubsection{Rich Context Components}

Our Rich Context captures:

\textbf{Schema Metadata:} Tables, columns, data types, primary/foreign keys, constraints.

\textbf{Index Information:} Index names, types (B-tree/Hash/GIN), columns, and selectivity metrics. For example:
\begin{verbatim}
{
  "idx_emp_dept": {
    "columns": ["department_id"],
    "type": "B-tree",
    "selectivity": 0.05,
    "guidance": "Use for department filtering"
  }
}
\end{verbatim}

\textbf{JOIN Path Graph:} We use BFS to compute shortest JOIN paths between any two tables via foreign key relationships, enabling correct multi-table query generation.

\textbf{Field Semantics:} Annotations for special fields (foreign keys, IDs, enumerations, timestamps) to aid LLM understanding.

\textbf{Database Dialect Features:} Syntax differences across databases:
\begin{itemize}
    \item String concatenation: \texttt{||} (PostgreSQL/SQLite) vs \texttt{CONCAT()} (MySQL)
    \item LIMIT syntax: \texttt{LIMIT n OFFSET m} vs \texttt{LIMIT m, n}
    \item Date arithmetic: \texttt{INTERVAL} vs \texttt{DATE\_ADD()}
\end{itemize}

\subsection{Runtime SQL Generation}

\subsubsection{Tool Design}

We employ a minimal tool set for capability decoupling:

\textbf{verify\_sql:} Validates SQL syntax via dry-run (\texttt{EXPLAIN}) without execution, providing fast feedback and execution plans.

\textbf{execute\_sql:} Executes SQL to verify correctness, retrieve missing information, or test performance.

\textbf{update\_rich\_context:} Updates context when information becomes stale or incorrect.

\subsubsection{ReAct Loop for SQL Generation}

\begin{algorithm}[h]
\caption{ReAct SQL Generation}
\begin{algorithmic}[1]
\STATE \textbf{Input:} Query $q$, Rich Context $C$
\STATE \textbf{Output:} SQL $s$
\STATE Initialize state $S_0 \leftarrow \{q, C\}$
\FOR{$t = 1$ to $T_{\max}$}
    \STATE $\theta_t \leftarrow \text{LLM}(S_{t-1})$ \COMMENT{Generate thought}
    \STATE $a_t \leftarrow \text{SelectAction}(\theta_t)$ \COMMENT{Choose tool}
    \IF{$a_t = \text{answer}(s)$}
        \RETURN $s$
    \ENDIF
    \STATE $o_t \leftarrow \text{Execute}(a_t)$ \COMMENT{Execute tool}
    \STATE $S_t \leftarrow S_{t-1} \cup \{\theta_t, a_t, o_t\}$
\ENDFOR
\end{algorithmic}
\end{algorithm}

\subsubsection{Database-Specific Syntax Guidance}

To handle SQL dialect differences, we inject database-specific syntax rules into prompts:

\begin{verbatim}
Database Type: SQLite
CRITICAL: Follow SQLite syntax strictly.
- Use || for string concatenation
- No LIMIT offset without LIMIT clause
- Date functions: date(), datetime(), strftime()
\end{verbatim}

This significantly reduces cross-database syntax errors.

\subsection{Dataset Calibration}

We systematically calibrated Spider dev through:
\begin{enumerate}
    \item \textbf{Error Identification:} Automated detection of syntax errors, execution failures, and schema mismatches
    \item \textbf{Manual Annotation:} Expert review of 221 problematic cases
    \item \textbf{Correction:} Fixed annotation errors, added \texttt{true\_question} for ambiguous queries, and corrected \texttt{true\_gold} SQL
    \item \textbf{Verification:} Validated corrections through execution testing
\end{enumerate}

This ensures reliable evaluation and fair comparison with existing methods.
